%==============================================================================
%  02-introduction.tex — Introduction & Research Scope
%==============================================================================
\strategyPart{02 \textbar\ Introduction}{From Traditional Work to AI-Based Organizations}

\section{Background of the study}

The rapid maturation of generative AI, predictive systems, assistants,
autonomous agents and multi-agent systems has moved artificial intelligence
from a laboratory capability to an operating-model decision. Organizations
everywhere are transitioning from traditional digital transformation toward
AI-enabled operating models. The evidence of adoption is overwhelming:

\begin{itemize}
  \item \textbf{88\% of organizations} now use AI in at least one business
        function — yet only \textbf{7\%} report that AI is fully scaled across
        the organization [McKinsey State of AI 2025].
  \item \textbf{78\% of organizations} reported using AI in at least one
        function in Stanford's 2024 industry survey, while global private AI
        investment reached roughly \textbf{\$252~billion} in 2024
        [Stanford AI Index 2025].
  \item \textbf{64\% of employees} report struggling with the time and energy
        to do their job — ``digital debt'' — and are 3.5$\times$ more likely
        to struggle with innovation [Microsoft Work Trend Index 2023].
  \item Generative AI alone could add \textbf{\$2.6--\$4.4~trillion} in annual
        value across 63 use cases [McKinsey 2023].
\end{itemize}

\section{The research problem}

There is a persistent gap between \textbf{AI experimentation} and
\textbf{meaningful organizational transformation}. Organizations purchase AI
tools and launch pilot projects, but they frequently fail to redesign
workflows, establish governance, train employees, or measure business value.
Deloitte's enterprise study — surveying \textbf{3,235 business and technology
leaders across 24 countries} — documents the transition from experimentation
to broader adoption as the central challenge of the current cycle
[Deloitte State of AI in the Enterprise 2026]. Gartner predicts that
\textbf{30\% of generative-AI projects will be abandoned after proof of
concept} through 2025 [Gartner 2023].

\section{Research purpose}

This survey develops a structured understanding of how traditional
organizations can transform into \textbf{AI-assisted}, \textbf{AI-enabled},
or \textbf{AI-native} organizations. It pursues four outcomes:

\begin{enumerate}
  \item Identify and compare the major organizational transformation models.
  \item Compare real-world success and failure case studies.
  \item Classify the major implementation challenges.
  \item Develop a practical, phased transformation framework.
\end{enumerate}

\section{Research objectives}

\begin{itemize}
  \item Define AI-based organizational transformation and its maturity dimensions.
  \item Compare established transformation frameworks (McKinsey, WEF, Deloitte,
        Stanford, NIST, OECD).
  \item Examine AI use across organizational functions and measure performance change.
  \item Analyze successful and unsuccessful case studies with documented results.
  \item Identify technical, organizational, human, ethical and regulatory challenges.
  \item Propose a transition roadmap from traditional work to AI-assisted work.
\end{itemize}

\section{Primary research question}

\begin{center}
  \yBanner[color=inNavy]{How can traditional organizations transform their operating models, processes, workforce and governance to achieve sustainable value from AI?}
\end{center}

Secondary questions include: What distinguishes AI \textit{adoption} from AI
\textit{transformation}? Which organizational capabilities are required to
scale AI? How does AI change jobs, skills and managerial responsibility? What
are the most common reasons AI initiatives fail? How can organizations combine
human expertise with assistants and autonomous agents? And how should AI
transformation success be measured?
